\section{The Architect's Question}\label{the-architects-question}

After this chapter we should be able to answer the question every
stakeholder actually asks --- \emph{what does this really cost?} ---
with a defensible number, not a guess. We will build a
total-cost-of-ownership model that separates capital, operating, and
opportunity costs, then price the canonical workload across three
delivery modes: self-hosted on our own cards, cloud GPU instances, and a
managed inference API. After this chapter, ``we save money by
self-hosting'' is either proven or refuted by arithmetic, and the
break-even point between modes is explicit.

\section{Concept}\label{concept}

Total Cost of Ownership (TCO) is the sum of everything required to
deliver a capability over its lifetime, not just the price tag of the
hardware or the per-token rate. It has three cost families:

\begin{enumerate}
\def\labelenumi{\arabic{enumi}.}
\item
  \textbf{Capital (capex)} --- the one-time cost of the cards, servers,
  networking, racks, and cooling that make up the system. For
  self-hosted inference this dominates up front: an 8×H100 server with
  networking and infrastructure carries a large price. {[}2° FACT{]}
\item
  \textbf{Operating (opex)} --- the recurring cost of running it:
  electricity and cooling, staff, software/licenses, cloud instance
  rental, and replacement/amortization of hardware. For cloud and
  managed API this is the whole cost; for self-hosted it compounds on
  capex.
\item
  \textbf{Opportunity} --- the cost of the \emph{alternative not taken}:
  the team time spent operating vs building, the delay to market, and
  the engineering hours that do not scale with either card or token
  counts. This is the most omitted and often the largest term.
\end{enumerate}

The architect's job is not to minimize capex, nor opex, nor provider
bill --- it is to minimize \textbf{the number that matters for the
organization's decision}, which usually means total cost per useful
request at the required quality and SLO.

\section{Mental Model}\label{mental-model}

Think of TCO as \textbf{a utilization-weighted average over the whole
system, not a sticker price.} Three mental moves keep the arithmetic
honest:

\begin{itemize}
\item
  \textbf{Amortize, don't stare at the price tag.} A \$400K server is an
  asset consumed over \textasciitilde5 years; spread it into
  dollars-per-month, then into dollars-per-request, before comparing to
  a per-token API rate. {[}2° DERIVED{]}
\item
  \textbf{Follow the token, not the GPU.} The thing being bought is
  completed, SLO-meeting requests. Two systems with the same card count
  can differ 10× in cost per request if one has 3× the goodput per card,
  better utilization, or no idle coasting.
\item
  \textbf{Include what does not scale.} Staff, power, networking
  engineering, and the latency of procurement are sunk or recurring
  whether utilization is 20\% or 80\%. A low-cost-per-card system that
  burns engineers can be the most expensive.
\end{itemize}

The decision shape is a \textbf{break-even graph}: at low volume,
managed/cloud-and-shut-down wins (no idle capex); at high volume,
self-hosted cards amortize below the per-token rate. The crossover point
is what we compute.

\section{Worked Example}\label{worked-example}

\subsection{Pricing the Canonical Workload Across Three
Modes}\label{pricing-the-canonical-workload-across-three-modes}

We price the canonical 70B enterprise-Q\&A workload at the canonical
traffic (\textasciitilde10 rps average, \textasciitilde40 rps peak,
\textasciitilde9,200 input + \textasciitilde300 output tokens per
request). {[}1P §14{]}

\textbf{Mode 1 --- self-hosted on owned cards (8×H100).}

\begin{itemize}
\tightlist
\item
  Capex: assume \textasciitilde\$350K for a fully built 8×H100 server
  (cards, host, NVSwitch, networking, rack, cooling), amortized over 5
  years → \textasciitilde\$70K/year, \textasciitilde\$5,800/month. {[}2°
  FACT, illustrative{]}
\item
  Opex: power + cooling for 8×H100 at \textasciitilde50 kW,
  \textasciitilde\$0.15/kWh, always on → \textasciitilde\$65K/year;
  staff/engineering alloc \textasciitilde\$120K/year; total opex
  \textasciitilde\$185K/year. {[}2° DERIVED, illustrative{]}
\item
  Monthly total ≈ \$5,800 (capex) + \$15,400 (opex) ≈
  \textbf{\$21,200/month}.
\item
  Per request at \textasciitilde10 rps average = \textasciitilde864,000
  requests/day ≈ 26M/month. \textbf{Cost ≈ \$21,200 / 26M ≈
  \$0.00082/request ≈ \$0.82 per 1,000 requests.} {[}2° DERIVED{]}
  \emph{(worked example, not reference)}
\end{itemize}

\textbf{Mode 2 --- cloud GPU instances (rent, shut down when idle).}

\begin{itemize}
\tightlist
\item
  An 8×H100 on-demand instance \textasciitilde\$3.50/hr (Ch4
  reconciled), paid only while running, say \textasciitilde40\% duty
  cycle to cover peaks → \textasciitilde\$1.40/hr average effective →
  \textasciitilde\$1,000/month. {[}2° DERIVED{]}
\item
  Per request: at 10 rps the same \textasciitilde26M requests/month.
  \textbf{Cost ≈ \$1,000 / 26M ≈ \$0.000038/request ≈ \$0.04 per 1,000
  requests.} {[}2° DERIVED{]} \emph{(cloud beats self-host per request
  ONLY at low duty cycle; the \$/hr must cover staff \& integration
  too.)}
\end{itemize}

\textbf{Mode 3 --- managed inference API.}

\begin{itemize}
\tightlist
\item
  A hosted frontier-70B-class API at \textasciitilde\$0.002/input +
  \textasciitilde\$0.008/output per 1K tokens (typical as of 2026). The
  per-request token cost is
\end{itemize}

\[
\text{cost/req} = p_\text{in} \cdot \frac{I}{1000} + p_\text{out} \cdot \frac{O}{1000} = 0.002 \times 9.2 + 0.008 \times 0.3 \approx \$0.0184 + \$0.0024 \approx \$0.0208 \approx \$20.8 \text{ per 1,000 requests}
\]

{[}2° FACT, illustrative{]} - At the canonical 25.9M requests/month:
\textasciitilde\$539,000/month --- an order of magnitude above
self-host. {[}2° DERIVED{]}

\textbf{Reading the result.} Self-hosted (\textasciitilde\$0.82/1K reqs)
beats the API (\textasciitilde\$20.8/1K reqs) by \textasciitilde25× on
pure token cost at this volume --- but only because we assume steady
near-canonical utilization plus in-house staff we are not separately
billing. Cloud-on-demand (\$0.04/1K) looks cheapest only if duty cycle
is low AND we ignore staff/integration. The honest TCO answer for a
business-critical, always-on service at canonical volume is
\textbf{self-hosted cards}, with the caveat that the staff term is the
real swing factor. {[}2° DERIVED{]}

\textbf{Parametric sensitivity is the durable form.} All the fixed
inputs above --- GPU price, electricity, staffing, capacity, API price
--- are {[}ILLUSTRATIVE{]} scenario values; the \emph{structure} of the
model is what generalizes. The architect can see the decision flip with
utilization and volume: if effective utilization is low
(\textasciitilde20\%, e.g.~a demo or dev workload), self-hosted cards
sit idle while depreciating, and managed on-demand wins; near the
crossover (\textasciitilde50--60\% utilization at canonical volume) the
two are close and the decision hinges on the staff/risk terms; only at
sustained high utilization (\textasciitilde80\%+) does self-hosted
clearly win on cost-per-request. Rather than memorize one answer, the
durable tool is a small parametric model --- exact same arithmetic with
GPU \$/hr, \$/kWh, staffing, and API price as inputs --- which the
architect re-runs with the customer's real numbers (\hyperref[chap:23]{Chapter 23}). This is
what makes TCO a decision \emph{framework}, not a single verdict. A
runnable, parameterised version of exactly this model ships with the
book as \texttt{render/tco\_calc.py} (defaults reproduce these numbers;
override GPU \$/hr, \$/kWh, staffing, and API price as flags).

\textbf{\hyperref[tab:16.1]{Table 16-1}} --- TCO comparison across delivery modes
(illustrative worked example)

\begin{longtable}[]{@{}
  >{\raggedright\arraybackslash}p{0.1667\linewidth}
  >{\raggedright\arraybackslash}p{0.1667\linewidth}
  >{\raggedright\arraybackslash}p{0.1667\linewidth}
  >{\raggedright\arraybackslash}p{0.1667\linewidth}
  >{\raggedright\arraybackslash}p{0.1667\linewidth}
  >{\raggedright\arraybackslash}p{0.1667\linewidth}@{}}
\toprule\noalign{}
\begin{minipage}[b]{\linewidth}\raggedright
Mode
\end{minipage} & \begin{minipage}[b]{\linewidth}\raggedright
Capex
\end{minipage} & \begin{minipage}[b]{\linewidth}\raggedright
Opex/month
\end{minipage} & \begin{minipage}[b]{\linewidth}\raggedright
Cost/1K req
\end{minipage} & \begin{minipage}[b]{\linewidth}\raggedright
\$/month @26M req
\end{minipage} & \begin{minipage}[b]{\linewidth}\raggedright
When it wins
\end{minipage} \\
\midrule\noalign{}
\endhead
\bottomrule\noalign{}
\endlastfoot
Self-hosted 8×H100 & \textasciitilde\$350K & \textasciitilde\$21.2K &
\textasciitilde\$0.82 & \textasciitilde\$21K & steady high utilization,
staff available \\
Cloud on-demand & \$0 & \textasciitilde\$1K (40\% duty) &
\textasciitilde\$0.04 & \textasciitilde\$1K+staff & low duty cycle,
bursty, no staff for ops \\
Managed API & \$0 & usage & \textasciitilde\$21 & \textasciitilde\$546K
& tiny volume, fastest time-to-value \\

\label{tab:16.1}\end{longtable}

\emph{(Figures are worked-example estimates as of Q4 2026, not vendor
quotes; treat as {[}2° DERIVED{]} illustrative, {[}VERIFY{]} before
budgeting.)}

\section{Measurement}\label{measurement}

Four metrics keep the TCO model honest:

\begin{enumerate}
\def\labelenumi{\arabic{enumi}.}
\tightlist
\item
  \textbf{Cost per useful (SLO-meeting) request} --- the decision
  number; divide total cost by goodput-meeting requests (Ch6), not raw
  tokens.
\item
  \textbf{Effective utilization} --- actual goodput ÷ theoretical
  capacity across the billing period; low utilization is the hidden tax
  in both capex and rented instances.
\item
  \textbf{Break-even volume} --- the requests/month where self-host
  cost-per-request crosses managed API; below it, don't buy cards.
\item
  \textbf{Staff/duty-cycle adjusters} --- the engineering-hours and
  duty-cycle assumptions that dominate the comparisons; measure them,
  don't assume.
\end{enumerate}

\section{Common Mistakes}\label{common-mistakes}

\begin{itemize}
\tightlist
\item
  \textbf{Comparing sticker price to per-token rate.} A \$/1K-token API
  quote already includes amortization; a card price tag does not.
  Amortize both.
\item
  \textbf{Ignoring the staff/ops term.} Self-host ``saves money'' on
  paper and loses the moment two engineers are consumed full-time.
\item
  \textbf{Assuming 100\% utilization.} Idle cards are still burning
  electricity and depreciation; bake duty cycle in.
\item
  \textbf{Forgetting opportunity cost.} A 6-month procurement cycle
  versus a same-week API integration is a real cost measured in
  time-to-value.
\item
  \textbf{Using peak capacity for cost-per-request.} Peak sizing (40
  rps) overstates cost if average (10 rps) utilization is what governs.
\end{itemize}

\section{Architecture Consequence}\label{architecture-consequence}

TCO is the final gate in the design loop (Ch12 candidates → Ch14
benchmark → Ch16 cost). It settles which candidate survives: the 70B
8×H100 self-host passes the capability screen, meets the SLO in the
benchmark, and wins on TCO at canonical volume --- completing the chain
from vague requirement to a defensible, costed architecture. It also
feeds fleet decisions (Ch17-18): the break-even curve is what justifies
buying a second host versus renting peaks, and the staff term is what
pushes toward consolidation (Ch18).

\begin{figure}
\centering
\pandocbounded{\includegraphics[keepaspectratio,alt={\hyperref[fig:16.1]{Fig 16.1} --- TCO break-even: self-hosted vs managed API {[}2° DERIVED{]}},width=\maxwidth{\textwidth},keepaspectratio]{figures/fig-16-1601.pdf}}
\caption{\hyperref[fig:16.1]{Fig 16.1} --- TCO break-even: self-hosted vs managed API {[}2°
DERIVED{]}}

\label{fig:16.1}\end{figure}

\emph{\hyperref[fig:16.1]{Fig 16.1} --- Break-even: monthly cost vs monthly volume for
self-hosted and managed API across three serve modes, crossing at the
\textasciitilde1.02 M-request scale.}

\section{What We Still Don't Know}\label{what-we-still-dont-know}

\begin{itemize}
\tightlist
\item
  \textbf{Real equipment/power/staff numbers} for a given deployment are
  {[}VERIFY{]} site-specific; the illustrative quotes above must be
  re-priced.
\item
  \textbf{How rapidly GPU list and amortization prices fall} over the
  5-year horizon is {[}HYPOTHESIS{]}; newer cards (H200/B200-class)
  change the per-card economics.
\item
  \textbf{The true staff overhead of self-hosting} (SRE time, security,
  upgrades) is {[}HYPOTHESIS{]} until the team bills its own time
  honestly.
\end{itemize}

\section{End-of-Chapter Mini-Case}\label{end-of-chapter-mini-case}

A startup's CTO shows the architect a warrant to buy two 8×H100 servers
for the new internal RAG assistant, citing ``we'll save on token fees.''
Applying TCO discipline, the architect does not dispute the raw token
math. Instead they price it: at this startup's \emph{actual} early
volume --- \textasciitilde5 rps average, \textasciitilde5M
requests/month, and only 0.5 engineers fully available for GPU ops ---
the break-even against a managed API sits just below that volume, and
the staff term dwarfs the token savings. The honest number: self-hosting
two servers \textasciitilde\$40K/month all-in vs
\textasciitilde\$105K/month API at current volume, but the two-server
option consumes a full engineer to keep utilization and reliability up
--- an opportunity cost the startup cannot yet afford. The architect's
verdict: \textbf{start on the managed API now, revisit self-host at
\textasciitilde2× this volume or when a second engineer frees up}, and
re-run the same model then. The CTO did not buy the servers --- the
break-even graph, with staff honestly included, made the answer
self-evident.
